\documentclass[10pt,a4paper]{amsart}
\usepackage{amsmath,amssymb,amsthm,graphicx}
\usepackage{url}
\usepackage{hyperref}
\usepackage{enumerate}

\DeclareMathOperator{\Cliff}{Clif\/f}
\DeclareMathOperator{\Gal}{Gal}
\DeclareMathOperator{\GL}{GL}
\DeclareMathOperator{\Inn}{Inn}
\DeclareMathOperator{\Mat}{Mat}
\DeclareMathOperator{\Out}{Out}
\DeclareMathOperator{\SL}{SL}
\DeclareMathOperator{\SO}{SO}
\DeclareMathOperator{\SU}{\mathbf{SU}}
\numberwithin{equation}{section}

\newtheorem{prop}{Proposition}[section]
\newtheorem{cor}[prop]{Corollary}
\newtheorem{lem}[prop]{Lemma}
{\theoremstyle{definition}
\newtheorem{defn}[prop]{Definition}
\newtheorem{notation}[prop]{Notation}
\newtheorem{rem}[prop]{Remark}
\newtheorem{ques}[prop]{Question}
\newtheorem{eg}[prop]{Example}
\newtheorem{erratum}[prop]{Erratum}}
\newcommand{\Res}[1]{\mathop{\mathbf{R}_{#1}}}

\newcommand{\ResOne}[1]{\mathop{\mathbf{R}^{(1)}_{#1}}}

\newcommand{\Lie}[1]{\mathfrak{\lowercase{#1}}}

\newcommand{\arXivNumber}{1410.2339}

\allowdisplaybreaks
\title{Tits-algebra compatibility criterion for real-universal semisimple rational Lie algebras}
\author{Dave Witte Morris}
\date{doi:10.3842/sigma.2015.034; CC BY 4.0}
\begin{document}
\maketitle
\noindent\textit{Editorial note.} The task counts the final semisimple direct-sum criterion only. The existence of some universal rational form(Proposition1.2, old item2007), separate classical-type classifications and the restriction-of-scalars proposition are not additional counted claims.\par\medskip
\noindent\textit{Editorial note.} In927, the second factor is printed with lambda\_i again. The preceding923–924 and following931–933 explicitly use lambda\_i and lambda\_j. Both real factors are H under the hypotheses, so this duplicated index does not change the calculation; preserve the original formula and disclose the index.\par\medskip
\noindent\textit{Editorial note.} The source explicitly writes in912 that only the forward direction is displayed and that the argument is reversible. The allowed-input boundary does not supply the converse as a theorem. The solver still has to check all global dominant weights: star-fixed weights give products of the common quaternion class, and non-star-fixed weights have center L. For the latter case, the Tits classes of an already R-universal outer factor split over its inner field L, using the retained root-lattice invariance and weight-center reasoning461–492. This is the scope of the author compressed reversibility claim; it is flagged for independent completeness review rather than adding an alternate authored proof.\par\medskip
\noindent\textit{Editorial note.} When comparing Tits algebras after tensor products, tilde denotes Brauer equivalence; D itself is the division representative. A product H tensor H is matrix algebra/Brauer trivial, not a division algebra. The source uses that distinction in921–933.\par\medskip
\noindent\textit{Editorial note.} The quaternionic type-D classification720–765, including the unresolved pair of real half-spin classes in759, is not an input or part of this task. Conditions(1) and the retained general representation criterion concern actual R-universal factors, not that case table.\par\medskip
\noindent\textit{Editorial note.} All of415–493 remains a local proof obligation. Only the external standard Tits results and basic Brauer/Galois facts listed below may be freely invoked.\par\medskip
\noindent\textit{Editorial scope.} Keep the entire author proof, including the explicitly compressed reverse direction. Separate notes expose its required case split and the duplicated index; no replacement proof or formula is inserted.\par\medskip
\noindent\textit{Allowed standard prerequisites (editorial).}\begin{itemize}
\item Tits highest-weight descent and Tits division algebra theorem. External classical classification of irreducible F-representations, the center as the fixed field of the star stabilizer, and the corresponding Brauer class. Do not give RunivIff or SameFixed for free.
\item Classical representation and Galois-cohomology foundations. Endomorphism algebras commute with scalar extension; dominant weights can be shifted into the root lattice; the star action has finite image. The current descent and common-field argument must be proved.
\item Brauer theory over Q, R and C. R division algebras with these centers are R/C/H; H tensor H is Brauer trivial, quaternion classes have order2, a central division algebra is determined by its Brauer class, and Tits classes add with central characters. No semisimple compatibility criterion is included as an input.
\end{itemize}\medskip
\begin{prop}
%\label{SS}
A~direct sum $\Lie G = \Lie G_1 \oplus \dots \oplus \Lie G_r$ of simple Lie algebras over~$\mathbb{Q}$ is
$\mathbb{R}$-universal if and only if:
\renewcommand{\labelenumi}{{\rm (\theenumi)}}
\begin{enumerate}\itemsep=0pt
\item
\label{SS-Gi}
each $\Lie G_i$ is $\mathbb{R}$-universal,
\item
\label{SS-star}
either $\Lie G$ is inner, or it becomes inner over an imaginary quadratic extension~$L$ of~$\mathbb{Q}$, and the
$*$-action of $\Gal(\mathbb{C}/\mathbb{R})$ is the same as the $*$-action of $\Gal(L/\mathbb{Q})$.
\item
\label{SS-Di=Dj}
we have $D_{\Lie G_i, \mathbb{Q}}(\lambda_i) \cong D_{\Lie G_j, \mathbb{Q}}(\lambda_j)$ whenever $\lambda_i$
and~$\lambda_j$ are $*$-invariant dominant weights of $G_i$ and~$G_j$, such that $D_{\Lie G_i, \mathbb{Q}}(\lambda_i)
\neq \mathbb{Q}$ and $D_{\Lie G_j, \mathbb{Q}}(\lambda_j) \neq \mathbb{Q}$, and
\item
\label{SS-Dsplit}
either $\Lie G$ is inner, or it becomes inner over an imaginary quadratic extension~$L$ of~$\mathbb{Q}$, such that
$D_{\Lie G_i, \mathbb{Q}}(\lambda_i)$ splits over~$L$, for every dominant weight~$\lambda_i$ of every $\Lie G_i$ that is
inner $($over~$\mathbb{Q})$.
\end{enumerate}
\end{prop}

\begin{defn}
Let $\Lie G_\mathbb{Q}$ be a~Lie algebra over~$\mathbb{Q}$.
(All Lie algebras and all representations are assumed to be f\/inite-dimensional.)
\begin{enumerate}\itemsep=0pt
\item $\Lie G_\mathbb{Q}$ is \emph{universal for real representations} (or \emph{$\mathbb{R}$-universal}, for short) if
every real representation of~$\Lie G_\mathbb{Q}$ has a~$\mathbb{Q}$-form~\cite[Def\/inition~7.1]{Morris-RRepsQForms}.
This means that if $\rho \colon \Lie G_\mathbb{Q} \to \Lie{GL}(n,\mathbb{R})$ is any ($\mathbb{Q}$-linear) Lie algebra
homomorphism, then there exists $M \in \GL(n,\mathbb{R})$, such that $M \rho(x) M^{-1} \in \Lie{GL}(n,\mathbb{Q})$, for
every $x \in \Lie G_\mathbb{Q}$.
\item $\Lie G_\mathbb{Q}$ is a~\emph{$\mathbb{Q}$-form} of a~real Lie algebra~$\Lie G_\mathbb{R}$ if $\Lie G_\mathbb{Q}
\otimes_\mathbb{Q} \mathbb{R} \cong \Lie G_\mathbb{R}$.
\end{enumerate}
\end{defn}

\section{Proof of the results stated in the introduction}\label{PfSect}

We begin by recalling a~result of J.~Tits that uses Galois cohomology to characterize the irreducible representations of
semisimple algebraic groups over f\/ields that are not algebraically closed.
However, we will state this work in the setting of Lie algebras, rather than algebraic groups.
(We deal only with semisimple Lie algebras and semisimple groups, and the f\/ields under consideration in this paper are
all of characteristic zero, so no dif\/f\/iculties arise in making this translation~\cite[Proposition~7.3.1(iii), p.~393]{DemazureGrothendieck-SGA3}.)

\begin{defn}[\protect{\cite[\S~4.2]{Tits-RepLinIrred}}] \label{BetaDefn}
Suppose $\Lie G$ is a~semisimple Lie algebra over a~subf\/ield~$F$ of~$\mathbb{C}$, and let $\mathbf{G}$ be the
corresponding simply connected, semisimple algebraic group over~$F$.
It is well known that there is a~(unique) quasi-split, simply connected algebraic group $\mathbf{G}^q$ over~$F$, and
a~$1$-cocycle $\xi \colon \Gal(\overline{F}/F) \to \overline{\mathbf{G}}{}^q$, where $\overline{\mathbf{G}}{}^q$ is the
adjoint group of~$\mathbf{G}^q$, such that $\mathbf{G}$ is~$F$-isomorphic to the Galois twist ${}^\xi \mathbf{G}^q$.
This cocycle determines a~cohomology class $[\xi] \in H^1(F; \overline{\mathbf{G}}{}^q)$.

Letting $Z(\mathbf{G}^q)$ be the center of $\mathbf{G}^q$, the short exact sequence $ e \to Z(\mathbf{G}^q) \to
\mathbf{G}^q \to \overline{\mathbf{G}}{}^q \to e $ yields a~corresponding long exact sequence of Galois cohomology sets,
including a~connecting map $\delta_* \colon H^1(F; \overline{\mathbf{G}}{}^q) \to H^2\bigl(F; Z(\mathbf{G}^q) \bigr)$.
Hence, we have a~cohomology class $\delta_* [\xi] \in H^2\bigl(F; Z(\mathbf{G}^q) \bigr)$.
(If we took a~bit more care to ensure that it is well-def\/ined, this would be the \emph{Tits class}
of~$\mathbf{G}$~\cite[p.~426]{KnusEtAl-BookOfInvols}.)

Now, f\/ix a~maximal~$F$-torus~$\mathbf{T}$ of~$\mathbf{G}^q$ that contains a~maximal~$F$-split torus, and
suppose~$\lambda$ is a~weight of~$\mathbf{T}$ that is invariant under the $*$-action of the Galois group $\Gal(\overline{F}/F)$.
Then the restriction of~$\lambda$ to $Z(\mathbf{G}^q)$ is a~$\Gal(\overline{F}/F)$-equivariant homomorphism from
$Z(\mathbf{G}^q)$ to the group~$\boldsymbol\mu$ of roots of unity in~$\mathbb{C}$, so it induces a~homomorphism
$\lambda_* \colon H^2\bigl(F; Z(\mathbf{G}^q) \bigr) \to H^2\bigl(F; \boldsymbol\mu \bigr)$.
Therefore, we may def\/ine
\begin{gather*}
\beta_{\Lie G,F}(\lambda) = \lambda_* \delta_* [\xi] \in H^2(F; \boldsymbol\mu).
\end{gather*}
\end{defn}

\begin{rem}
For a~semisimple group~$\mathbf{G}$ that is def\/ined over a~f\/ield~$F$, the def\/inition of the $*$-action of
$\Gal(\overline{F}/F)$ on the weights of a~maximal torus~$\mathbf{T}$ presupposes that $\mathbf{T}$ is def\/ined over~$F$
and contains a~maximal~$F$-split torus of~$\mathbf{G}$~\cite[p.~66]{PlatonovRapinchukBook}.
When $\mathbf{G}$ is def\/ined over~$\mathbb{Q}$, we will use the $*$-actions of both $\Gal(\overline{\mathbb{Q}}/\mathbb{Q})$ and $\Gal(\mathbb{C}/\mathbb{R})$, so we will assume that $\mathbf{T}$ is def\/ined
over~$\mathbb{Q}$ (so it is also def\/ined over~$\mathbb{R}$) and contains both a~maximal $\mathbb{Q}$-split torus and
a~maximal $\mathbb{R}$-split torus.
To see that such a~choice of~$\mathbf{T}$ is always possible, let $\mathbf{T}_1$ be any maximal $\mathbb{Q}$-split torus
of~$\mathbf{G}$.
Then there is a~maximal $\mathbb{Q}$-torus~$\mathbf{T}$ of the centralizer $\mathbf{C}_{\mathbf{G}}(\mathbf{T}_1)$ that
contains a~maximal $\mathbb{R}$-split torus~\cite[Corollary~2 of Proposition~7.8, p.~418]{PlatonovRapinchukBook}, and it is clear
that $\mathbf{T}$ has the desired properties.
\end{rem}

\begin{prop}[\protect{\cite[Theorem~7.2 and Lemma~7.4]{Tits-RepLinIrred}}]\label{HighWt}
Suppose $\Lie G$ is a~semisimple Lie algebra over a~subfield~$F$ of~$\mathbb{C}$, and~$\lambda$ is a~dominant weight.
Then:
\renewcommand{\labelenumi}{{\rm \theenumi.}}
\begin{enumerate}\itemsep=0pt
\item There is an irreducible representation ${}^F \rho_\lambda \colon \Lie G \to \Lie{GL}(n, F)$, for some~$n$, such
that \mbox{${}^F \rho_\lambda \otimes_F \mathbb{C}$} has an irreducible summand with highest weight~$\lambda$.
Furthermore, ${}^F \rho_\lambda$ is unique up to isomorphism.

\item ${}^F \rho_{\lambda_1} \cong {}^F \rho_{\lambda_2}$ if and only if~$\lambda_1$ and~$\lambda_2$ are in the same
orbit of the $*$-action of~$\Gal(\overline{F}/F)$.

\item
\label{HighWt-IrredIff}
${}^F \rho_\lambda \otimes_F \mathbb{C}$ is irreducible if and only if:
\begin{enumerate}[{\rm (a)}]\itemsep=0pt
\item
%\label{HighWt-IrredIff-Action}
$\lambda$ is invariant under the $*$-action of~$\Gal(\overline{F}/F)$, and
\item
%\label{HighWt-IrredIff-coho}
$\beta_{\Lie G,F}(\lambda)$ is the trivial element of $H^2(F; \boldsymbol\mu)$.
\end{enumerate}
\end{enumerate}
Furthermore, every~$F$-irreducible representation of~$\Lie G$ is isomorphic to ${}^F \rho_{\lambda}$, for some
dominant weight~$\lambda$.
\end{prop}

\begin{notation}
We assume $\Lie G$ is a~semisimple Lie algebra over a~f\/ield~$F$ of characteristic~$0$, that~$\lambda$~is a~dominant
weight, and the other notation of Def\/inition~\ref{BetaDefn}.
Furthermore, we let $Z^*$ be the f\/inite, abelian group of all homomorphisms from $Z(\mathbf{G}^q)$ to~$\boldsymbol\mu$.
\end{notation}

\begin{prop}[\protect{\cite[Corollary~3.5, \S~4.2, and Lemma~7.4]{Tits-RepLinIrred}}]\label{Beta=D}
Let~$L$ be the center of the algebra $D_{\Lie G, F}(\lambda):= \mathrm{End}_{\Lie G}({}^F \rho_\lambda)$ $($which,
by Schur's lemma, is a~division algebra$)$.
Then:
\renewcommand{\labelenumi}{{\rm (\theenumi)}}
\begin{enumerate}\itemsep=0pt
\item
\label{Beta=D-invt}
$\Gal(\overline{F}/L) = \{ \sigma \in \Gal(\overline{F}/F) \mid \sigma(\lambda) = \lambda \}$ $($for the $*$-action
of $\Gal(\overline{F}/F))$, and
\item $\beta_{\Lie G, L}(\lambda) = [D_{\Lie G, F}(\lambda)]$, after identifying $H^2(L; \boldsymbol\mu)$ with the
Brauer group of~$L$.
\end{enumerate}
\end{prop}

\begin{defn}
We will call $D_{\Lie G, F}(\lambda)$ the \emph{Tits algebra} of~$\Lie G$ corresponding to the weight~$\lambda$
over~$F$.
(However, this name is used in the literature for a~slightly dif\/ferent algebra that is Brauer equivalent to $D_{\Lie G,
F}(\lambda)$~\cite[\S~27A, p.~377]{KnusEtAl-BookOfInvols}.)
\end{defn}


\begin{lem}
\label{RunivIff}
If $\Lie G$ is a~semisimple Lie algebra over~$\mathbb{Q}$, then the following are equivalent:
\renewcommand{\labelenumi}{{\rm \theenumi.}}
\begin{enumerate}\itemsep=0pt
\item
\label{RunivIff-univ}
$\Lie G$ is $\mathbb{R}$-universal.
\item
\label{RunivIff-irred}
${}^\mathbb{Q} \rho_\lambda \otimes_\mathbb{Q} \mathbb{R} \cong {}^\mathbb{R} \rho_\lambda$, for every dominant
weight~$\lambda$.
\item
\label{RunivIff-DivAlg}
$D_{\Lie G, \mathbb{Q}}(\lambda) \otimes_\mathbb{Q} \mathbb{R}$ is a~division algebra, for every dominant
weight~$\lambda$.
\item
\label{RunivIff-D}
$D_{\Lie G, \mathbb{Q}}(\lambda) \otimes_\mathbb{Q} \mathbb{R} \cong D_{\Lie G, \mathbb{R}}(\lambda)$, for every
dominant weight~$\lambda$.
\end{enumerate}
\end{lem}

\begin{proof}
($\ref{RunivIff-univ} \Rightarrow \ref{RunivIff-irred}$) Let~$\rho$ be a~$\mathbb{Q}$-form of ${}^\mathbb{R}
\rho_\lambda$.
Then~$\rho$ must be $\mathbb{Q}$-irreducible (since $\rho \otimes_\mathbb{Q} \mathbb{R} \cong {}^\mathbb{R}
\rho_\lambda$ is $\mathbb{R}$-irreducible) and
\begin{gather*}
\rho \otimes_\mathbb{Q} \mathbb{C} \cong (\rho \otimes_\mathbb{Q} \mathbb{R}) \otimes_\mathbb{R} \mathbb{C} \cong
{}^\mathbb{R} \rho_\lambda \otimes_\mathbb{R} \mathbb{C}
\end{gather*}
has an irreducible summand with highest weight~$\lambda$.
Therefore $\rho \cong {}^\mathbb{Q} \rho_\lambda$.

($\ref{RunivIff-irred} \Rightarrow \ref{RunivIff-univ}$) ${}^\mathbb{Q} \rho_\lambda$ is a~$\mathbb{Q}$-form of
${}^\mathbb{R} \rho_\lambda$.

($\ref{RunivIff-irred} \Leftrightarrow \ref{RunivIff-DivAlg} \Leftrightarrow \ref{RunivIff-D}$) We have
$\mathrm{End}_{\Lie G}({}^\mathbb{Q} \rho_\lambda \otimes_\mathbb{Q} \mathbb{R}) = D_{\Lie G, \mathbb{Q}}(\lambda)
\otimes_\mathbb{Q} \mathbb{R}$.
So ${}^\mathbb{Q} \rho_\lambda \otimes_\mathbb{Q} \mathbb{R}$ is irreducible (and hence equal to~${}^\mathbb{R}
\rho_\lambda$) if and only if $D_{\Lie G, \mathbb{Q}}(\lambda) \otimes_\mathbb{Q} \mathbb{R}$ is a~division algebra (and
hence equal to~$D_{\Lie G, \mathbb{R}}(\lambda)$).
\end{proof}

\begin{rem}
\label{OnlyZStar}
It is immediate from Def\/inition~\ref{BetaDefn} that if the $*$-invariant weights $\lambda_1$ and $\lambda_2$ have the same
restriction to $Z(\mathbf{G}^q)$ (that is, if they represent the same element of~$Z^*$), then $\beta_{\Lie G,
F}(\lambda_1) = \beta_{\Lie G, F}(\lambda_2)$.
Therefore, if~$\lambda$ is a~dominant weight in the root lattice (so~$\lambda$ is trivial on $Z(\mathbf{G}^q)$),
and~$\lambda$ is invariant under the $*$-action of $\Gal(\overline{F}/F)$, then $D_{\Lie G, F}(\lambda) = F$.
\end{rem}

\begin{lem}
\label{SameFixed}
Let  $\Lie G$ be a~semisimple Lie algebra over~$\mathbb{Q}$.
If~$\Lie G$ is $\mathbb{R}$-universal, then $\Lie G$ is inner over some imaginary quadratic extension~$L$
of~$\mathbb{Q}$, and the $*$-action of $\Gal(L/\mathbb{Q})$ is the same as the $*$-action of~$\Gal(\mathbb{C}/\mathbb{R})$.
\end{lem}

\begin{proof}
Assume $\Lie G$ is not inner, for otherwise the desired conclusions are obvious.
Let~$L$ be the (unique) minimal extension of~$\mathbb{Q}$ over which $\Lie G$ becomes inner, and choose some dominant
weight~$\lambda$ that is not f\/ixed by any nontrivial element of~$\Gal(L/\mathbb{Q})$.
We know that~$L$ is the center of~$D_{\Lie G, \mathbb{Q}}(\lambda)$ (see Proposition~\ref{Beta=D}(\ref{Beta=D-invt})) and that $D_{\Lie
G, \mathbb{Q}}(\lambda) \otimes_\mathbb{Q} \mathbb{R}$ is a~division algebra (see Lemma~\ref{RunivIff}(\ref{RunivIff-DivAlg})).
Since~$L$ is not $\mathbb{Q}$, this implies it is an imaginary quadratic extension of~$\mathbb{Q}$.

It now suf\/f\/ices to show that every every $\Gal(\mathbb{C}/\mathbb{R})$-invariant dominant weight~$\lambda$ is also
$\Gal(L/\mathbb{Q})$-invariant.
Suppose not.
By replacing~$\lambda$ with an appropriate positive integer multiple, we may assume it is in the root lattice, so
$D_{\Lie G,\mathbb{R}}(\lambda) = \mathbb{R}$ (see Remark~\ref{OnlyZStar}).
Therefore, we must have $D_{\Lie G,\mathbb{Q}}(\lambda) = \mathbb{Q}$ (see Lemma~\ref{RunivIff}(\ref{RunivIff-D})).
However, this contradicts Proposition~\ref{Beta=D}(\ref{Beta=D-invt}), since~$\lambda$ is not
$\Gal(\overline{\mathbb{Q}}/\mathbb{Q})$-invariant.
\end{proof}

\begin{proof}
We prove only ($\Rightarrow$), but the argument is reversible.
(We use $\sim$ to denote ``is Brauer equivalent to''.)

(\ref{SS-Gi}) Any representation $\rho_i$ of~$\Lie G_i$ extends to a~representation of $\Lie G$ (that is $0$ on the
other simple factors), so it is clear that $\Lie G_i$ must be $\mathbb{R}$-universal.

(\ref{SS-star}) See Lemma~\ref{SameFixed}.
(This means that all $\Lie G_i$ become inner over the same quadratic extension~$L$.)

(\ref{SS-Di=Dj}) We have
\begin{gather*}
D_{\Lie G, \mathbb{R}}(\lambda_i + \lambda_j) \sim D_{\Lie G, \mathbb{R}}(\lambda_i) \otimes_\mathbb{R} D_{\Lie G,
\mathbb{R}}(\lambda_j) \\
\hphantom{D_{\Lie G, \mathbb{R}}(\lambda_i + \lambda_j)}{}
\cong \bigl(D_{\Lie G, \mathbb{Q}}(\lambda_i) \otimes_\mathbb{Q} \mathbb{R} \bigr)
\otimes_\mathbb{R} \bigl(D_{\Lie G, \mathbb{Q}}(\lambda_i) \otimes_\mathbb{Q} \mathbb{R} \bigr) \cong \mathbb{H}
\otimes_\mathbb{R} \mathbb{H} \sim \mathbb{R}.
\end{gather*}
Since $\Lie G$ is $\mathbb{R}$-universal, this implies $D_{\Lie G, \mathbb{Q}}(\lambda_i + \lambda_j) \sim\mathbb{Q}$,
from which we conclude that $D_{\Lie G, \mathbb{Q}}(\lambda_i) \otimes_\mathbb{Q} D_{\Lie G, \mathbb{Q}}(\lambda_j) \sim
\mathbb{Q}$, so $D_{\Lie G, \mathbb{Q}}(\lambda_i) \cong D_{\Lie G, \mathbb{Q}}(\lambda_j)$ (since these are quaternion
algebras).

(\ref{SS-Dsplit}) Let $\lambda_j$ be a~dominant weight of some (outer) simple factor $G_j$ that is not f\/ixed by the
$*$-action.
Then $\lambda_i + \lambda_j$ is not f\/ixed by the $*$-action, so
\begin{gather*}
D_{\Lie G, \mathbb{Q}}(\lambda_i + \lambda_j) = D_{\Lie G, L}(\lambda_i + \lambda_j) \sim D_{\Lie G_i, L}(\lambda_i)
\otimes_L D_{\Lie G_j, L}(\lambda_j) = D_{\Lie G_i, L}(\lambda_i) \otimes_L L \cong D_{\Lie G_i, L}(\lambda_i).
\end{gather*}
Since $D_{\Lie G, \mathbb{R}}(\lambda_i + \lambda_j) = D_{\Lie G, \mathbb{C}}(\lambda_i + \lambda_j) = \mathbb{C}$, and
$\Lie G$ is $\mathbb{R}$-universal, we conclude from Lemma~\ref{RunivIff}(\ref{RunivIff-D}) that $D_{\Lie G_i, L}(\lambda_i) = L$.
Since $D_{\Lie G, \mathbb{Q}}(\lambda_i) \otimes_\mathbb{Q} L \sim D_{\Lie G, L}(\lambda_i)$, this means that $D_{\Lie
G, \mathbb{Q}}(\lambda_i)$ splits over~$L$.
\end{proof}

\begin{thebibliography}{99}
\footnotesize\itemsep=0pt

\bibitem{DemazureGrothendieck-SGA3}
Demazure M., Grothendieck A., Sch\'emas en groups.~III. Structure des sch\'emas
  en groupes reductifs, S\'eminaire de g\'eom\'etrie alg\'ebrique du {B}ois
  {M}arie, 1962/1964, available at
  \url{http://library.msri.org/books/sga/sga/pdf/sga3-3.pdf}.

\bibitem{Harder-BerichtNeuere}
Harder G., Bericht \"uber neuere {R}esultate der {G}aloiskohomologie
  halbeinfacher {G}ruppen, \textit{Jber. Deutsch. Math.-Verein.} \textbf{70}
  (1968), 182--216, available at
  \url{http://www.digizeitschriften.de/en/dms/img/?PPN=PPN37721857X_0070&DMDID=dmdlog17}.

\bibitem{Kneser-LectGalCohoClassical}
Kneser M., Lectures on {G}alois cohomology of classical groups, \textit{Tata
  Institute of Fundamental Research Lectures on Mathematics}, Vol.~47, Tata
  Institute of Fundamental Research, Bombay, 1969, available at
  \url{http://www.math.tifr.res.in/~publ/ln/tifr47.pdf}.

\bibitem{KnusEtAl-BookOfInvols}
Knus M.A., Merkurjev A., Rost M., Tignol J.P., The book of involutions,
  \textit{American Mathematical Society Colloquium Publications}, Vol.~44,
  Amer. Math. Soc., Providence, RI, 1998.

\bibitem{Lam-IntroQuadForms}
Lam T.Y., Introduction to quadratic forms over f\/ields, \textit{Graduate Studies
  in Mathematics}, Vol.~67, Amer. Math. Soc., Providence, RI, 2005.

\bibitem{Morris-RRepsQForms}
Morris D., Real representations of semisimple {L}ie algebras have {${\mathbb
  Q}$}-forms, in Algebraic Groups and Arithmetic, Tata Institute of Fundamental
  Research, Mumbai, 2004, 469--490.

\bibitem{PlatonovRapinchukBook}
Platonov V., Rapinchuk A., Algebraic groups and number theory, \textit{Pure and
  Applied Mathematics}, Vol.~139, Academic Press, Inc., Boston, MA, 1994.

\bibitem{PrasadRapinchuk-Prescribed}
Prasad G., Rapinchuk A., On the existence of isotropic forms of semi-simple
  algebraic groups over number f\/ields with prescribed local behavior,
  \href{http://dx.doi.org/10.1016/j.aim.2006.01.001}{\textit{Adv. Math.}} \textbf{207} (2006), 646--660.

\bibitem{Raghunathan-ArithLattsSSGrps}
Raghunathan M.S., Arithmetic lattices in semisimple groups, \href{http://dx.doi.org/10.1007/BF02967980}{\textit{Proc.
  Indian Acad. Sci. Math. Sci.}} \textbf{91} (1982), 133--138.

\bibitem{Tits-Classification}
Tits J., Classif\/ication of algebraic semisimple groups, in Algebraic {G}roups
  and {D}iscontinuous {S}ubgroups ({P}roc. {S}ympos. {P}ure {M}ath., {B}oulder,
  {C}olo., 1965), Amer. Math. Soc., Providence, R.I., 1966, 33--62.

\bibitem{Tits-RepLinIrred}
Tits J., Repr\'esentations lin\'eaires irr\'eductibles d'un groupe r\'eductif
  sur un corps quelconque, \href{http://dx.doi.org/10.1515/crll.1971.247.196}{\textit{J.~Reine Angew. Math.}} \textbf{247} (1971),
  196--220.

\bibitem{Tits-StronglyInner}
Tits J., Strongly inner anisotropic forms of simple algebraic groups,
  \href{http://dx.doi.org/10.1016/0021-8693(90)90201-X}{\textit{J.~Algebra}} \textbf{131} (1990), 648--677.

\end{thebibliography}
\end{document}
